\chapter{The eALS Algorithm and its PySpark Implementation}
\label{ch:methods}

This chapter has two goals, matching the two things the assignment asks
for separately. The first is to explain the algorithm eALS actually is:
where its update rule comes from and why it is fast. The second is to walk
through how that algorithm was mapped onto Spark, with the actual code that
does it. All source files referenced here live under \texttt{src/} in the
project repository; the full listings are reproduced in
Appendix~\ref{app:code}.

\section{Deriving the element-wise update}
\label{sec:ealsdesign}

The paper's key move is to stop updating a user's latent vector $p_u$ as a
whole and instead optimize one coordinate $p_{uf}$ at a time, holding every
other coordinate fixed: ordinary coordinate descent, applied to
\eqref{eq:objective}. Setting $\partial L/\partial p_{uf}=0$ gives a
closed-form update for a single coordinate that needs no matrix inversion:
\begin{equation}
  \label{eq:puf_naive}
  p_{uf} = \frac{\sum_{i=1}^{N}(r_{ui}-\hat r_{ui}^f)\,w_{ui}\,q_{if}}
                {\sum_{i=1}^{N} w_{ui}\,q_{if}^2 + \lambda},
\end{equation}
where $\hat r_{ui}^f = \hat r_{ui} - p_{uf}q_{if}$ is the prediction with
factor $f$'s contribution removed. Taken at face value this is no better
than before: the sum still runs over all $N$ items, most of which $u$ never
touched. The paper's actual contribution is a reformulation that separates
the observed part of the sum (small, $|\mathcal{R}_u|$ terms) from the
missing part (large, $N-|\mathcal{R}_u|$ terms) and pre-computes the
missing part once per iteration instead of once per user. Define the
$K\times K$ cache
\[
  S^q = \sum_{i=1}^N c_i\, q_i q_i^\top ,
\]
which does not depend on $u$ at all. Because it does not depend on $u$, it
can be computed once (a single matrix multiplication, $O(NK^2)$) and reused
for \emph{every} user's update in that sweep. With $S^q$ available,
\eqref{eq:puf_naive} becomes
\begin{equation}
  \label{eq:puf}
  p_{uf} =
  \frac{\sum_{i\in\mathcal{R}_u}\big[w_{ui}r_{ui}-(w_{ui}-c_i)\hat r_{ui}^f\big]q_{if}
        - \sum_{k\neq f} p_{uk}\, s^q_{kf}}
       {\sum_{i\in\mathcal{R}_u}(w_{ui}-c_i)\,q_{if}^2 + s^q_{ff} + \lambda},
\end{equation}
where $s^q_{kf}$ denotes the $(k,f)$ entry of the matrix $S^q$. This is the
paper's Equation 12; the code comments in Appendix~\ref{app:code} cite it
by that number rather than by the label used here. Every sum in
Eq.~\ref{eq:puf} now runs only over $u$'s own observed items:
$O(|\mathcal{R}_u|)$ work instead of $O(N)$. The item update $q_{if}$ is
derived symmetrically, with one difference that deserves a flag, because it
is the easiest place to introduce a bug when implementing this: the
missing-data confidence $c_i$
belongs to the \emph{item}, not the user, so it appears as a constant
scalar multiplying the whole cache term, rather than as a per-neighbour
weight the way it does in \eqref{eq:puf}:
\begin{equation}
  \label{eq:qif}
  q_{if} =
  \frac{\sum_{u\in\mathcal{R}_i}\big[w_{ui}r_{ui}-(w_{ui}-c_i)\hat r_{ui}^f\big]p_{uf}
        - c_i\sum_{k\neq f} q_{ik}\, s^p_{kf}}
       {\sum_{u\in\mathcal{R}_i}(w_{ui}-c_i)\,p_{uf}^2 + c_i\, s^p_{ff} + \lambda},
\end{equation}
where $s^p_{kf}$ denotes the $(k,f)$ entry of $S^p = P^\top P$ (unweighted,
since popularity weighting in this formulation only lives on the item
side); this is the paper's Equation 13, again cited by that number in the
code comments. Getting Eq.~\ref{eq:qif}'s
placement of $c_i$ wrong, treating it as a per-neighbour term like in
Eq.~\ref{eq:puf} instead of a scalar multiplying the whole cache
contribution, was in fact the first bug caught while implementing this
project. Section~\ref{sec:correctness} describes how it was found.

\medskip
One more detail matters for both performance and correctness:
$\hat r_{ui}^f$ never needs to be recomputed from scratch. Because the
factors are updated one at a time in a fixed order, the current prediction
$\hat r_{ui}$ can be maintained incrementally: subtract $p_{uf}q_{if}$
before updating $p_{uf}$, then add the new $p_{uf}q_{if}$ back afterward,
turning an $O(K)$ dot product into an $O(1)$ update. This is what brings the total
per-iteration cost down from $O(MNK)$ (a naive coordinate-descent
implementation, still touching every entry) to
$O((M+N)K^2+|\mathcal{R}|K)$: linear in the number of observed
interactions rather than quadratic in the size of the whole matrix.
Algorithm~\ref{alg:eals} summarizes the full procedure.

\begin{algorithm}[H]
\caption{eALS -- one training iteration}
\label{alg:eals}
\begin{algorithmic}[1]
\Require $\mathcal{R}$, $K$, $\lambda$, item confidences $c$, current $P, Q$
\Ensure updated $P, Q$
\State $S^q \leftarrow \sum_{i=1}^N c_i\, q_i q_i^\top$
\ForAll{users $u = 1 \ldots M$ \textbf{(independent, parallel)}}
    \ForAll{factors $f = 1 \ldots K$}
        \State update $p_{uf}$ via Eq.~\ref{eq:puf}, using the cached $\hat r_{ui}$
        \State refresh $\hat r_{ui}$ for $i \in \mathcal{R}_u$ in $O(1)$ per item
    \EndFor
\EndFor
\State $S^p \leftarrow P^\top P$
\ForAll{items $i = 1 \ldots N$ \textbf{(independent, parallel)}}
    \ForAll{factors $f = 1 \ldots K$}
        \State update $q_{if}$ via Eq.~\ref{eq:qif}, using the cached $\hat r_{ui}$
        \State refresh $\hat r_{ui}$ for $u \in \mathcal{R}_i$ in $O(1)$ per user
    \EndFor
\EndFor
\end{algorithmic}
\end{algorithm}

One more piece of the algorithm is worth deriving before moving to the
implementation, since it is what produces every loss curve in this report:
checking convergence needs the value of the objective $L$ itself, not just
the updated $P$ and $Q$. Computed directly from \eqref{eq:objective}, this
would cost $O(MNK)$, dominated by the same missing-data sum that made the
naive update slow in the first place. The paper's Equation 14 avoids this
with the same trick used for the updates: reusing $S^q$,
\[
  \sum_{u=1}^{M}\sum_{i\notin\mathcal{R}_u} c_i\,\hat r_{ui}^{\,2}
  = \sum_{u=1}^{M} p_u^\top S^q p_u - \sum_{(u,i)\in\mathcal{R}} c_i\,\hat r_{ui}^{\,2},
\]
which turns the missing-data term of $L$ into a sum over just $M$ users
plus a correction over the (much smaller) observed set $\mathcal{R}$,
bringing the cost down to $O(|\mathcal{R}|+MK^2)$.
\texttt{EALS.compute\_loss} in \texttt{src/eals.py} implements exactly this
expression; it is what produced the loss values reported in
Section~\ref{sec:correctness} and plotted in the demonstration notebook
(Appendix~\ref{app:notebook}).

\section{Why this is embarrassingly parallel}
\label{sec:parallel}

The paper observes, but does not build, the following fact: within the
user-update loop of Algorithm~\ref{alg:eals}, $Q$ and $S^q$ are read-only
for the whole duration of the loop, and each user's update writes only to
that user's own row of $P$. No user's update reads any other user's
current or past state. The order in which users are processed therefore
does not affect the result, and the loop can be split across an arbitrary
number of workers with no coordination beyond having $Q$ and $S^q$
available to all of them. This is a stronger property than
``parallelizable with some loss of accuracy'', which is what one gets from
naively parallelizing SGD, where different workers' updates can conflict.
eALS's parallel execution is instead bit-for-bit equivalent, up to
floating-point summation order, to running it on a single machine. The same
argument applies to the item loop.

\section{Mapping the algorithm onto Spark}
\label{sec:spark}

Section~\ref{sec:parallel} argued that Algorithm~\ref{alg:eals}'s two
sweeps can be split across workers with no coordination and no loss of
accuracy. This section describes how that argument was turned into actual
Spark code, then checked for correctness.

\subsection{Design}

The implementation, \texttt{src/eals.py}, keeps $P$ and $Q$ as ordinary
NumPy arrays on the Spark driver. At the scale evaluated in this report
($33{,}326$ users, $21{,}901$ items, $K$ up to $128$), both matrices
comfortably fit in memory, and keeping them there avoids the overhead of shuffling
factor updates through Spark's DataFrame machinery on every single
coordinate. What \emph{is} distributed is the actual computation described
in Section~\ref{sec:parallel}: the id space $\{0,\dots,M-1\}$ (respectively
$\{0,\dots,N-1\}$) is partitioned with \texttt{sc.parallelize(...).glom()},
$Q$, $S^q$ and each user's interaction list are broadcast to every
executor, and the closed-form update of Eq.~\ref{eq:puf} runs as ordinary
NumPy code inside each partition, with no further communication until the
results are collected back to the driver and written into $P$. Listing
\ref{lst:update_users} shows the actual per-partition update routine.

\begin{lstlisting}[language=Python, caption={Closed-form user update, Eq.~\ref{eq:puf}, executed inside one Spark partition (\texttt{src/eals.py}).}, label={lst:update_users}]
def _update_users_chunk(
    ids_chunk, user_ratings, P, Q, item_conf, Sq, reg, w_obs,
):
    """Eq. (12): closed-form update of p_uf for every user in
    `ids_chunk`, iterating f=1..K."""
    K = P.shape[1]
    out = np.empty((len(ids_chunk), K), dtype=np.float64)
    for row, u in enumerate(ids_chunk):
        neighbors = user_ratings.get(u)
        pu = P[u].copy()
        if neighbors is None or len(neighbors) == 0:
            out[row] = pu
            continue
        Qu = Q[neighbors]                # |R_u| x K
        cu = item_conf[neighbors]        # |R_u|, c_i per neighbor item
        rhat = Qu @ pu
        for f in range(K):
            rhat_f = rhat - pu[f] * Qu[:, f]
            numer = np.sum((w_obs - (w_obs - cu) * rhat_f) * Qu[:, f]) \
                    - (pu @ Sq[:, f] - pu[f] * Sq[f, f])
            denom = np.sum((w_obs - cu) * Qu[:, f] ** 2) + Sq[f, f] + reg
            new_val = numer / denom if denom > 1e-12 else 0.0
            rhat = rhat_f + new_val * Qu[:, f]   # O(1) prediction refresh
            pu[f] = new_val
        out[row] = pu
    return out
\end{lstlisting}

Every identifier in Listing~\ref{lst:update_users} corresponds to a
quantity in Eq.~\ref{eq:puf}: \texttt{reg} is $\lambda$, \texttt{w\_obs} is
$w_{ui}$ for observed entries, \texttt{cu} holds $c_i$ for each item in
$u$'s neighbor list, and \texttt{Sq} is the cached matrix $S^q$.
\texttt{rhat\_f} is $\hat r_{ui}^f$, the first summand of
\texttt{numer} is the observed-data term, the subtraction is
$\sum_{k\neq f} p_{uk}s^q_{kf}$ (computed as a full dot product minus the
$k=f$ term, which is cheaper than an explicit loop excluding one index),
and \texttt{denom} is the closed-form denominator. The reassignment
\texttt{rhat = rhat\_f + new\_val * Qu[:, f]}, executed right after
\texttt{new\_val} is computed and before the loop moves to the next
factor, is the $O(1)$ prediction maintenance from
Section~\ref{sec:ealsdesign}: it folds the newly computed $p_{uf}$ back
into the running prediction before moving on to factor $f+1$, matching the
Gauss-Seidel semantics of the paper's Algorithm~1 (each factor update
immediately sees the previous factors' new values, within the same user,
in the same sweep).

\medskip
The driver-side training loop, shown in abbreviated form in
Listing~\ref{lst:fit}, follows Algorithm~\ref{alg:eals} step for step: it
recomputes $S^q$ from the current $Q$, dispatches the parallel user sweep,
waits for it to complete and folds the results back into $P$, then repeats
the same pattern for $S^p$ and the item sweep.

\begin{lstlisting}[language=Python, caption={Driver-side training loop (excerpt from \texttt{EALS.fit}, \texttt{src/eals.py}).}, label={lst:fit}]
for it in range(cfg.n_iters):
    # ---- update user factors: Algorithm 1, lines 4-11 -------------
    Sq = factor_cache(self.Q, self.item_conf)
    Q_b, Sq_b, P_local = sc.broadcast(self.Q), sc.broadcast(Sq), self.P

    def upd_users(ids_chunk, ur=ur_b, Qb=Q_b, Sqb=Sq_b, icb=ic_b,
                  cfg=cfg, P=P_local):
        ids_chunk = np.asarray(ids_chunk)
        res = _update_users_chunk(ids_chunk, ur.value, P, Qb.value,
                                   icb.value, Sqb.value, cfg.reg, cfg.w_obs)
        return list(zip(ids_chunk.tolist(), res.tolist()))

    for uid, vec in user_chunks.flatMap(upd_users).collect():
        self.P[uid] = vec
    Q_b.unpersist(); Sq_b.unpersist()

    # ---- update item factors: Algorithm 1, lines 12-19 -------------
    Sp = factor_cache(self.P)
    # ... symmetric flatMap over item_chunks, updating self.Q ...
\end{lstlisting}

\medskip
Two implementation choices deserve to be made explicit, since both were
deliberate and both trade a bit of theoretical elegance for something that
actually runs reliably.

\begin{bulletlist}
  \item \textbf{Broadcast, not shuffle.} An alternative design would keep
        $P$ and $Q$ as Spark DataFrames and express the update as a join
        between users and their rated items. This was tried early on and
        abandoned: it forces a shuffle on every single coordinate update
        (there are $K$ of them per sweep), and Spark's DataFrame/Catalyst
        overhead per job dwarfs the few milliseconds of actual NumPy work
        each partition does at this problem size. Broadcasting the (small)
        $K\times K$ cache and the (moderate) $Q$ matrix, and only
        distributing the large but embarrassingly parallel per-user
        computation, keeps the number of Spark jobs at exactly two per
        iteration instead of $2K$.
  \item \textbf{One partition, one NumPy loop.} Each partition still
        contains a plain Python \texttt{for} loop over its assigned users.
        This is the one place a purely mathematical description of the
        algorithm and its actual runtime cost part ways: for the dataset
        used in this report, the inner loop is dominated by Python function
        call overhead rather than the O($K$)-per-item arithmetic itself,
        which is the reason Section~\ref{sec:scale_k} finds close-to-linear
        rather than perfectly linear scaling in $K$.
\end{bulletlist}

\subsection{A note on correctness}
\label{sec:correctness}

Because the two update rules, Eq.~\ref{eq:puf} and Eq.~\ref{eq:qif}, look
almost identical but are not, the implementation was checked in two
independent ways before trusting it on real data.

\medskip
First, a synthetic dataset with a known block structure was constructed:
three disjoint communities of 15 users and 15 items each, where every
interaction stays within its own community. A correct factorization has no
reason to ever mix latent directions across communities, so a trained model
should rank a user's own community's items far above the other two. This
is exactly what happens: after 20 iterations the training loss decreases
monotonically from 145.7 to 83.8, and the top-10 recommended items for a
probed user are all drawn from that user's own community (see
\texttt{tests/test\_eals\_synthetic.py} and the corresponding notebook
cell, Appendix~\ref{app:notebook}). Monotonic loss decrease is not a given
here, since coordinate descent on a non-convex objective can still increase
the loss if a single update is wrong, so this is a meaningful check and not
a formality.

\medskip
Second, and more usefully, an early version of \texttt{\_update\_items\_chunk}
reused the same helper function as the user update, passing an array of
all-ones as the per-neighbor confidence instead of correctly scaling the
whole cache term by the item's own $c_i$ as Eq.~\ref{eq:qif} requires. This
silently zeroed out the $(w_{ui}-c_i)$ correction on the item side, without
crashing or producing an obviously wrong-looking loss curve. It is the
kind of bug that a synthetic sanity check alone would not necessarily have
caught, since the community-separation test still passes with weaker item
updates. It was only found by re-deriving Eq.~\ref{eq:qif} term by term
against the implementation, which is why Section~\ref{sec:ealsdesign}
above spells out the derivation in enough detail to make that comparison
possible for a reader, rather than just stating the final formula.
